% Source: https://artofproblemsolving.com/community/c3h1955915_geo_help
% Stable identity: aops-topic-1955915-q01
Let the largest circle have radius $R$. A circle of radius $r$ tangent to both sides at a corner has its centre on the diagonal, at distance $r\sqrt2$ from the corner. So the three centres lie on the diagonal (figure shows the corner, magnified).

\begin{asy}
usepackage("amsmath");
size(6cm);
defaultpen(fontsize(10pt));
real R=1, a=0.1715728752538097, s=0.029437251522859347;
draw((0,0)--(0.6,0),linewidth(0.9));draw((0,0)--(0,0.6),linewidth(0.9));
draw(arc((1,1),1,192,258),linewidth(0.8));
draw(circle((a,a),a),blue+linewidth(0.8));
draw(circle((s,s),s),red+linewidth(0.8));
draw((0,0)--(0.6,0.6),dashed+gray(0.5));
dot((a,a),blue);dot((s,s),red);
label("big circle (radius $R$)",(0.47,0.76),fontsize(8pt));
label("$r_1$",(a,a),E,blue+fontsize(8pt));
label("$r_2$",(s,0),(1.2,-1.6),red+fontsize(8pt));
label("\mbox{}",(min(currentpicture,user=true)+max(currentpicture,user=true))/2,black); // keeps the text after the figure black when LaTeX typesets the labels inline
\end{asy}

The big circle's centre is at distance $R\sqrt2$ from the corner, and its nearest point on the diagonal at distance $R\sqrt2-R$. The medium circle (radius $r_1$) touches it from the corner side, so
\[
r_1\sqrt2+r_1=R\sqrt2-R\quad\Longrightarrow\quad r_1=R\cdot\dfrac{\sqrt2-1}{\sqrt2+1}=(3-2\sqrt2)R .
\]
The same argument with the medium circle in place of the big one gives $r_2=(3-2\sqrt2)r_1=(3-2\sqrt2)^2R$. The ratio of areas is
\[
\left(\dfrac{r_2}{R}\right)^2=(3-2\sqrt2)^4=(17-12\sqrt2)^2=\boxed{577-408\sqrt2}\approx0.00075 .
\]
